% ============================================================================
%  Rapport technique — Système de bases de données réparties « EShop »
%
%  Compilation (pdfLaTeX, deux passes pour la table des matières) :
%      pdflatex rapport.tex
%      pdflatex rapport.tex
%
%  Paquets requis (TeX Live « full » ou installation à la carte) :
%      inputenc, fontenc, lmodern, babel (french), geometry, xcolor,
%      listings, graphicx, booktabs, enumitem, microtype, fancyhdr,
%      tikz (+ arrows.meta), hyperref
%
%  Sur Debian/Ubuntu :  sudo apt install texlive-full
%  ou, minimal :        texlive-latex-recommended texlive-latex-extra
%                       texlive-lang-french texlive-pictures
% ============================================================================
\documentclass[11pt,a4paper]{article}

% --- Encodage et police -----------------------------------------------------
\usepackage[utf8]{inputenc}
\usepackage[T1]{fontenc}
\usepackage{lmodern}
\usepackage[french]{babel}
\usepackage{microtype}
\usepackage{amsmath}
\usepackage{amssymb}

% --- Mise en page -----------------------------------------------------------
\usepackage[a4paper,margin=2.4cm]{geometry}
\usepackage{graphicx}
\usepackage{booktabs}
\usepackage{enumitem}
\usepackage{xcolor}
\usepackage{fancyhdr}

% --- Schémas ----------------------------------------------------------------
\usepackage{tikz}
\usetikzlibrary{arrows.meta}

% --- Captures d'écran (emplacements réservés) -------------------------------
% Utilisation :  \screenshot{captures/mon-image.png}{légende de remplacement}
% Si le fichier existe, il est inséré ; sinon un cadre indiquant le chemin
% attendu s'affiche, ce qui permet de compiler le rapport avant d'avoir pris
% les captures. Déposez les images dans docs/captures/ puis recompilez.
\newcommand{\screenshot}[2]{%
  \IfFileExists{#1}{%
    \includegraphics[width=0.95\linewidth]{#1}%
  }{%
    \fbox{\begin{minipage}[c][5cm][c]{0.9\linewidth}%
      \centering\itshape\color{cmt}%
      [\,capture à insérer\,]\\[6pt]
      Fichier attendu : \texttt{#1}\\[6pt]
      #2%
    \end{minipage}}%
  }%
}

% --- Couleurs du thème ------------------------------------------------------
\definecolor{codebg}{RGB}{248,249,250}
\definecolor{codeframe}{RGB}{222,226,230}
\definecolor{kw}{RGB}{0,92,197}
\definecolor{cmt}{RGB}{106,115,125}
\definecolor{str}{RGB}{163,21,21}
\definecolor{num}{RGB}{0,108,156}
\definecolor{accent}{RGB}{0,86,179}
\definecolor{glob}{RGB}{219,234,254}
\definecolor{site}{RGB}{255,237,213}
\definecolor{bkp}{RGB}{220,242,228}
\definecolor{dash}{RGB}{236,238,240}

% --- Extraits de code (listings) -------------------------------------------
\usepackage{listings}
\renewcommand{\lstlistingname}{Extrait}
\renewcommand{\lstlistlistingname}{Liste des extraits de code}

\lstdefinestyle{base}{
  backgroundcolor=\color{codebg},
  basicstyle=\ttfamily\footnotesize,
  keywordstyle=\color{kw}\bfseries,
  commentstyle=\color{cmt}\itshape,
  stringstyle=\color{str},
  numberstyle=\tiny\color{cmt},
  numbers=left,
  numbersep=9pt,
  frame=single,
  rulecolor=\color{codeframe},
  framesep=6pt,
  xleftmargin=16pt,
  xrightmargin=4pt,
  breaklines=true,
  breakatwhitespace=true,
  showstringspaces=false,
  tabsize=2,
  captionpos=b,
  columns=fullflexible,
  keepspaces=true,
  literate=%
    {é}{{\'e}}1 {è}{{\`e}}1 {ê}{{\^e}}1 {ë}{{\"e}}1
    {à}{{\`a}}1 {â}{{\^a}}1 {ä}{{\"a}}1
    {î}{{\^i}}1 {ï}{{\"i}}1
    {ô}{{\^o}}1 {ö}{{\"o}}1
    {ù}{{\`u}}1 {û}{{\^u}}1 {ü}{{\"u}}1
    {ç}{{\c c}}1 {œ}{{\oe}}1 {Œ}{{\OE}}1
    {É}{{\'E}}1 {È}{{\`E}}1 {Ê}{{\^E}}1 {À}{{\`A}}1 {Ç}{{\c C}}1
    {×}{{$\times$}}1 {≥}{{$\geq$}}1 {≤}{{$\leq$}}1
    {→}{{$\rightarrow$}}1 {∧}{{$\wedge$}}1 {σ}{{$\sigma$}}1
    {’}{{'}}1 {‘}{{'}}1 {“}{{"}}1 {”}{{"}}1
    {–}{{--}}1 {—}{{---}}1 {…}{{\ldots}}1
}

\lstdefinestyle{sql}{
  style=base,
  language=SQL,
  morekeywords={REPLACE,PROCEDURE,TRIGGER,PRAGMA,ROWNUM,EXECUTE,
    IMMEDIATE,EXCEPTION,RAISE,NVL,ROUND,COMMIT,ROLLBACK,EACH,BEFORE,
    AFTER,LINK,SYNONYM,SEQUENCE,PLUGGABLE,DATABASE,CONTAINER,SESSION,
    USING,IDENTIFIED,CONNECT,EXPLAIN,PLAN,TRUNCATE,BOOLEAN,IS,LOOP},
}

\lstdefinestyle{ts}{
  style=base,
  morekeywords={const,let,var,async,await,export,import,from,type,
    interface,return,try,catch,finally,new,function,as,of,Promise,
    void,if,else,for,while,extends,implements,public,private,readonly},
  sensitive=true,
  morecomment=[l]{//},
  morecomment=[s]{/*}{*/},
  morestring=[b]',
  morestring=[b]",
  morestring=[b]`,
}

\lstdefinestyle{yaml}{
  style=base,
  morecomment=[l]{\#},
  morestring=[b]",
  morestring=[b]',
}

% --- En-têtes et pieds de page ---------------------------------------------
\pagestyle{fancy}
\fancyhf{}
\fancyhead[L]{\small\itshape Bases de données réparties : EShop}
\fancyhead[R]{\small\thepage}
\fancyfoot[C]{\small Rapport technique}
\renewcommand{\headrulewidth}{0.4pt}
\renewcommand{\footrulewidth}{0.0pt}

% --- Liens hypertexte (chargé en dernier) ----------------------------------
\usepackage[colorlinks=true,linkcolor=accent,urlcolor=accent,
  citecolor=accent,linktoc=all]{hyperref}

% ============================================================================
\begin{document}

% --------------------------- PAGE DE GARDE ----------------------------------
\begin{titlepage}
\centering
\vspace*{1.5cm}

{\Large\scshape Rapport technique}\\[0.4cm]
{\large Module : Bases de données réparties}\\[2.0cm]

\rule{\linewidth}{0.6pt}\\[0.5cm]
{\huge\bfseries Conception et mise en œuvre\\[0.3cm]
d'un système de bases de données réparties}\\[0.4cm]
{\Large Fragmentation horizontale, synchronisation\\ et optimisation sous Oracle~23}\\[0.5cm]
\rule{\linewidth}{0.6pt}\\[2.0cm]

{\large\textbf{Application : plateforme de commerce « EShop »}}\\[2.5cm]

\begin{tabular}{rl}
\textbf{Auteur} & Ilyass Bougati \\[0.3cm]
\textbf{Technologies} & Oracle Database 23 Free, Docker Compose, \\
                      & PL/SQL, Next.js, node-oracledb \\[0.3cm]
\textbf{Date} & \today \\
\end{tabular}

\vfill
{\small Document généré automatiquement à partir du code source du projet.}
\end{titlepage}

% --------------------------- TABLE DES MATIÈRES -----------------------------
\tableofcontents
\newpage

% ============================================================================
\section{Introduction et objectifs}
% ============================================================================

Ce rapport présente la conception et la réalisation d'un \textbf{système de
gestion de bases de données réparties} appliqué à une plateforme de commerce
électronique fictive nommée \emph{EShop}. Le projet illustre les concepts
fondamentaux des bases de données distribuées : la \emph{fragmentation
horizontale} des données sur plusieurs sites, la \emph{transparence de
localisation} pour l'application cliente, la \emph{synchronisation}
automatique entre un nœud central et des nœuds satellites, et enfin
l'\emph{optimisation} des requêtes locales et distribuées.

\subsection{Contexte fonctionnel}

Le schéma applicatif modélise une boutique en ligne classique : des clients
passent des commandes, chaque commande étant détaillée en lignes de commande
qui référencent des produits, eux-mêmes classés par catégorie et rattachés à
des fournisseurs. La table volumineuse et au cœur de l'activité est
\texttt{LIGNECOMMANDES} : c'est précisément cette table que l'on cherche à
\emph{fragmenter} et à \emph{répartir} sur plusieurs serveurs.

\subsection{Objectifs techniques}

Le cahier des charges se décline en six questions (Q1 à Q6), toutes couvertes
par l'implémentation :

\begin{description}[leftmargin=2.2cm,style=nextline]
  \item[\textbf{Q1}] Définir le schéma global et fragmenter horizontalement la
    table \texttt{LIGNECOMMANDES} selon des prédicats métier.
  \item[\textbf{Q2}] Doter chaque site de procédures stockées d'insertion, de
    suppression et de mise à jour sur son fragment local.
  \item[\textbf{Q3}] Garantir l'intégrité référentielle de chaque fragment
    (clés primaires, clés étrangères, rapatriement des lignes parentes).
  \item[\textbf{Q4}] Synchroniser automatiquement les sites depuis le nœud
    global au moyen de déclencheurs (\emph{triggers}).
  \item[\textbf{Q5}] Analyser et optimiser une requête analytique au moyen de
    plans d'exécution (\texttt{EXPLAIN PLAN}) et d'index.
  \item[\textbf{Q6}] Reconstruire une vue globale (chiffre d'affaires par
    catégorie) à partir des fragments distants via des liens de base de
    données.
\end{description}

\subsection{Pile technologique}

L'ensemble repose sur \textbf{Oracle Database 23 Free} (image
\texttt{gvenzl/oracle-free}), chaque nœud étant un conteneur \textbf{Docker}
orchestré par \textbf{Docker Compose}. La logique métier est écrite en
\textbf{PL/SQL}. Un \textbf{tableau de bord web} (Next.js + \texttt{node-oracledb})
complète le dispositif pour visualiser en temps réel l'état des nœuds et le
routage des données.

% ============================================================================
\section{Architecture générale}
% ============================================================================

\subsection{Vue d'ensemble}

L'architecture suit un modèle \textbf{hub-and-spoke} (étoile) : un nœud
\emph{global} central détient le jeu de données complet et joue le rôle
d'unique point d'écriture ; autour de lui gravitent des nœuds \emph{site} qui
n'hébergent qu'un fragment des données, ainsi qu'un nœud de \emph{sauvegarde}.
Tous les conteneurs communiquent sur un réseau Docker privé,
\texttt{maroc-db-net}, et se résolvent mutuellement par des liens de base de
données Oracle (\emph{database links}).

\begin{figure}[h]
\centering
\begin{tikzpicture}[
    font=\small,
    box/.style={draw,thick,rounded corners,align=center,
                minimum width=3.1cm,minimum height=1.25cm},
    glob/.style={box,fill=glob},
    site/.style={box,fill=site},
    bkp/.style={box,fill=bkp},
    dash/.style={box,fill=dash,minimum width=4cm},
    lk/.style={-{Latex[length=2.2mm]},thick},
    lkd/.style={-{Latex[length=2.2mm]},thick,dashed},
]
% Réseau Docker (rectangle de fond, dessiné avant les noeuds)
\draw[dashed,rounded corners,draw=accent!60]
      (-6.9,-3.7) rectangle (6.9,3.6);
\node[accent,anchor=north west] at (-6.8,3.5) {\footnotesize\bfseries maroc-db-net};

% Noeuds base de données
\node[glob] (g)  at (0,1.7)   {\textbf{oracle-global}\\[1pt]\scriptsize G\_PDB, port 1521\\[1pt]\scriptsize jeu complet + triggers};
\node[site] (s1) at (-4.4,-1.4){\textbf{oracle-site-1}\\[1pt]\scriptsize S1\_PDB, port 1522\\[1pt]\scriptsize fragment R1};
\node[site] (s2) at (4.4,-1.4) {\textbf{oracle-site-2}\\[1pt]\scriptsize S2\_PDB, port 1523\\[1pt]\scriptsize fragment R2};
\node[bkp]  (b)  at (0,-2.6)   {\textbf{oracle-backup}\\[1pt]\scriptsize B\_PDB, port 1524\\[1pt]\scriptsize copie + job horaire};

% Application
\node[dash] (app) at (0,5.0) {\textbf{Tableau de bord (Next.js)}\\[1pt]\scriptsize node-oracledb (pool de connexions)};

% Triggers : global -> sites
\draw[lk] (g.west)  to[bend right=18] node[above,sloped,font=\scriptsize]{triggers} (s1.north);
\draw[lk] (g.east)  to[bend left=18]  node[above,sloped,font=\scriptsize]{triggers} (s2.north);
% DB links : sites -> global (synonymes)
\draw[lk] (s1.east) to[bend left=10]  node[below,font=\scriptsize]{DB link} (b.west);
\draw[lk] (s2.west) to[bend right=10] node[below,font=\scriptsize]{DB link} (b.east);
% Backup -> global (rechargement horaire)
\draw[lk] (b.north) -- node[right,font=\scriptsize]{pull horaire} (g.south);

% Application -> noeuds
\draw[lk]  (app.south) -- node[right,font=\scriptsize]{écritures} (g.north);
\draw[lkd] (app.west)  to[bend right=30] node[left,font=\scriptsize]{lectures} (s1.west);
\draw[lkd] (app.east)  to[bend left=30]  node[right,font=\scriptsize]{lectures} (s2.east);
\end{tikzpicture}
\caption{Architecture répartie en étoile sur le réseau Docker privé. Les
écritures arrivent toujours sur le nœud global ; les déclencheurs propagent
chaque ligne vers le site propriétaire ; le nœud de sauvegarde recharge
périodiquement l'intégralité du jeu de données.}
\label{fig:archi}
\end{figure}

\subsection{Les nœuds du système}

Chaque nœud est une base de données enfichable (\emph{Pluggable Database}, PDB)
hébergée dans son propre conteneur. Le tableau~\ref{tab:noeuds} récapitule la
topologie pour le scénario~1.

\begin{table}[h]
\centering
\begin{tabular}{@{}llll@{}}
\toprule
\textbf{Nœud} & \textbf{Conteneur} & \textbf{Port hôte} & \textbf{Rôle} \\
\midrule
Global    & \texttt{oracle-global} & 1521 & Jeu de données complet, déclencheurs de synchro \\
Site 1    & \texttt{oracle-site-1} & 1522 & Fragment R1 \\
Site 2    & \texttt{oracle-site-2} & 1523 & Fragment R2 \\
Sauvegarde& \texttt{oracle-backup} & 1524 & Copie complète, rechargement horaire \\
\bottomrule
\end{tabular}
\caption{Topologie des nœuds (scénario~1). Le scénario~2 est déployé sur une
pile parallèle aux ports 1531--1533.}
\label{tab:noeuds}
\end{table}

\subsection{Création des bases enfichables (PDB)}
\label{sec:pdb}

L'image \texttt{gvenzl/oracle-free:23} fournit une base \emph{multilocataire}
(\emph{multitenant}) : un conteneur racine \texttt{CDB\$ROOT} pouvant héberger
plusieurs bases enfichables (\emph{Pluggable Databases}). Plutôt que d'utiliser
la PDB livrée par défaut, \textbf{chaque nœud crée la sienne} au premier
démarrage (\texttt{G\_PDB}, \texttt{S1\_PDB}, \texttt{S2\_PDB},
\texttt{B\_PDB}). Cela donne à chaque base un nom de service explicite (ce qui
rend les alias TNS et les liens lisibles) et isole complètement les jeux de
données les uns des autres.

La procédure est identique sur les quatre nœuds (seul le nom change). On se
place dans le conteneur racine, on \emph{clone} la PDB modèle \texttt{pdbseed}
en recopiant ses fichiers via \texttt{FILE\_NAME\_CONVERT}, on déclare un
administrateur local \texttt{pdb\_admin}, puis on ouvre la PDB et on
\textbf{sauvegarde son état} afin qu'elle se rouvre automatiquement à chaque
redémarrage du conteneur.

\begin{lstlisting}[style=sql,caption={Création d'une base enfichable depuis \texttt{CDB\$ROOT} (extrait de \texttt{01\_create\_schema.sql}, nœud global).},label={lst:pdb}]
-- 1) Se placer dans le conteneur racine de la base multilocataire
ALTER SESSION SET CONTAINER = CDB$ROOT;

-- 2) Cloner la PDB modele (pdbseed) en une nouvelle base enfichable
CREATE PLUGGABLE DATABASE G_PDB
  ADMIN USER pdb_admin IDENTIFIED BY Admin123
  ROLES = (DBA)
  DEFAULT TABLESPACE USERS
    DATAFILE '/opt/oracle/oradata/FREE/G_PDB/users01.dbf'
    SIZE 250M AUTOEXTEND ON
  FILE_NAME_CONVERT = (
    '/opt/oracle/oradata/FREE/pdbseed/',   -- source : la PDB modele
    '/opt/oracle/oradata/FREE/G_PDB/'      -- cible  : les fichiers de G_PDB
  );

-- 3) Ouvrir la PDB et l'enregistrer aupres de l'auditeur (listener)
ALTER PLUGGABLE DATABASE G_PDB OPEN;
ALTER SYSTEM REGISTER;

-- 4) Conserver l'etat OPEN au redemarrage du conteneur
ALTER PLUGGABLE DATABASE G_PDB SAVE STATE;

-- 5) Basculer dans la nouvelle PDB pour y creer le schema
ALTER SESSION SET CONTAINER = G_PDB;
\end{lstlisting}

Deux instructions méritent attention. \texttt{ALTER SYSTEM REGISTER} force
l'enregistrement immédiat du service auprès de l'\emph{auditeur}, sans quoi un
lien de base de données entrant pourrait ne pas résoudre le service pendant la
première minute. \texttt{SAVE STATE} évite d'avoir à rouvrir manuellement la PDB
(elle est \texttt{MOUNTED} mais fermée par défaut après un redémarrage).

\subsection{Orchestration par Docker Compose}

Le fichier \texttt{compose.yaml} décrit les conteneurs, leurs volumes
persistants, leurs limites de ressources, et surtout l'\textbf{ordre de
démarrage}. Les scripts SQL placés dans \texttt{scripts-*/} sont montés dans
\texttt{/container-entrypoint-initdb.d/} et exécutés automatiquement au premier
démarrage de chaque conteneur.

Le point délicat est la \textbf{dépendance de démarrage} : les sites ne doivent
s'initialiser qu'une fois le nœud global prêt, car leurs scripts rapatrient des
données depuis celui-ci. Cela est garanti par la combinaison d'un
\texttt{healthcheck} sur le global et d'une clause
\texttt{depends\_on:\ condition:\ service\_healthy} sur les sites. Le
\texttt{healthcheck} interroge une table sentinelle \texttt{INIT\_COMPLETE},
créée tout à la fin du dernier script du global : tant que la sentinelle est
absente, le global est considéré comme « non sain » et les sites attendent.

\begin{lstlisting}[style=yaml,caption={Service du nœud global : healthcheck par sentinelle (extrait de \texttt{compose.yaml}).},label={lst:compose}]
db-global:
  image: gvenzl/oracle-free:23-slim-faststart
  container_name: oracle-global
  ports:
    - "1521:1521"
  volumes:
    - global-data:/opt/oracle/oradata
    - ./scripts-global:/container-entrypoint-initdb.d/
    - ./tnsnames.ora:/opt/oracle/network/admin/tnsnames.ora
  healthcheck:
    test:
      [ "CMD-SHELL",
        "echo 'SELECT sentinel FROM pdb_admin.INIT_COMPLETE WHERE ROWNUM=1;' \
         | sqlplus -s s1_user/mon_mdp@//localhost:1521/G_PDB \
         | grep -q READY || exit 1" ]
    interval: 30s
    timeout: 10s
    retries: 20
    start_period: 120s
  networks:
    - maroc-db-net

db-site-1:
  # ...
  depends_on:
    db-global:
      condition: service_healthy   # le site attend que le global soit prêt
\end{lstlisting}

\subsection{Résolution de noms et liens de base de données}

Un fichier \texttt{tnsnames.ora} unique est partagé entre tous les conteneurs.
Il définit des alias réseau (\texttt{GLOBAL\_DB\_ALIAS}, \texttt{S1\_ALIAS},
\texttt{S2\_ALIAS}, \texttt{BACKUP\_ALIAS}) qui pointent vers les noms d'hôtes
Docker. Ces alias sont ensuite utilisés pour créer les \emph{database links},
qui sont le mécanisme par lequel un nœud Oracle accède aux objets d'un autre
nœud, comme s'ils étaient locaux.

\begin{lstlisting}[style=sql,caption={Alias TNS et création d'un lien de base de données.},label={lst:tns}]
-- tnsnames.ora : alias resolu via le DNS interne de Docker
GLOBAL_DB_ALIAS =
  (DESCRIPTION =
    (ADDRESS = (PROTOCOL = TCP)(HOST = db-global)(PORT = 1521))
    (CONNECT_DATA = (SERVER = DEDICATED)(SERVICE_NAME = G_PDB)))

-- Depuis un site, on cree un lien vers le global ...
CREATE DATABASE LINK link_to_global
    CONNECT TO s1_user IDENTIFIED BY mon_mdp
    USING 'GLOBAL_DB_ALIAS';

-- ... et l'on masque la localisation derriere des synonymes :
-- une requete sur PRODUITS interroge en realite le noeud global.
CREATE OR REPLACE SYNONYM PRODUITS FOR pdb_admin.PRODUITS@link_to_global;
\end{lstlisting}

L'usage de \textbf{synonymes} par-dessus les liens réalise la
\emph{transparence de localisation} : le code des procédures d'un site écrit
\texttt{SELECT \ldots\ FROM PRODUITS} sans savoir que la table physique réside
sur le nœud global.

\subsection{Utilisateurs, habilitations et communication inter-nœuds}
\label{sec:users}

\subsubsection*{Les comptes de chaque nœud}

Chaque PDB possède un propriétaire de schéma, \texttt{pdb\_admin}, créé comme
\texttt{ADMIN USER} à la naissance de la base (section~\ref{sec:pdb}). C'est le
seul compte puissant, et il n'est utilisé qu'au moment de l'initialisation par
les scripts montés dans le conteneur ; le système \emph{en fonctionnement} ne
s'en sert jamais à travers le réseau. Toute la communication inter-nœuds passe
par des \textbf{comptes de service dédiés}, au plus faible privilège possible
(tableau~\ref{tab:users}).

\begin{table}[h]
\centering
\begin{tabular}{@{}llll@{}}
\toprule
\textbf{Compte} & \textbf{Réside sur} & \textbf{Utilisé par} & \textbf{Privilèges accordés} \\
\midrule
\texttt{pdb\_admin} & toutes les PDB & scripts d'init & propriétaire du schéma (DBA local) \\
\texttt{s1\_user}   & \texttt{G\_PDB} & lien du Site 1 & \texttt{SELECT} (7 tables), \texttt{CREATE SESSION}, \\
                    &                 &                & \texttt{EXECUTE insert\_ligne\_commande} \\
\texttt{s2\_user}   & \texttt{G\_PDB} & lien du Site 2 & \texttt{SELECT} (7 tables), \texttt{CREATE SESSION} \\
\texttt{b\_user}    & \texttt{G\_PDB} & lien du Backup & \texttt{SELECT} (7 tables), \texttt{CREATE SESSION} \\
\texttt{g\_user}    & \texttt{S1\_PDB}, \texttt{S2\_PDB} & liens du global & \texttt{EXECUTE} (3 procédures), \texttt{SELECT} \\
                    &                 &                & (tables de fragment), \texttt{CREATE SESSION} \\
\bottomrule
\end{tabular}
\caption{Comptes fonctionnels et leurs privilèges. Aucun ne dispose des droits
\texttt{DBA}, \texttt{CREATE TABLE} ou d'écriture directe.}
\label{tab:users}
\end{table}

\subsubsection*{Le principe du moindre privilège}

Chaque compte de service ne reçoit que les droits strictement nécessaires à sa
mission, et rien de plus. Les comptes qui ne font que \emph{lire} le nœud
global (\texttt{s1\_user}, \texttt{s2\_user}, \texttt{b\_user}) n'obtiennent que
\texttt{SELECT} sur les tables et \texttt{CREATE SESSION} pour ouvrir une
connexion : ils ne peuvent ni modifier les données, ni créer d'objet, ni voir
les autres schémas. À l'inverse, le compte \texttt{g\_user}, par lequel le
global \emph{écrit} sur les sites, ne reçoit pas l'accès direct aux tables : il
n'a que le droit d'\texttt{EXECUTE} sur les trois procédures de routage, qui
encapsulent et contrôlent toutes les écritures.

\begin{lstlisting}[style=sql,caption={Habilitations au plus juste : aucun \texttt{GRANT ALL} (extraits de \texttt{04\_create\_roles.sql} et des scripts de site).},label={lst:grants}]
-- Sur le noeud GLOBAL : un compte de lien ne recoit que la LECTURE
GRANT SELECT ON pdb_admin.CLIENTS        TO s1_user;
GRANT SELECT ON pdb_admin.PRODUITS       TO s1_user;
GRANT SELECT ON pdb_admin.LIGNECOMMANDES TO s1_user;   -- ... idem pour les 7 tables
GRANT CREATE SESSION                     TO s1_user;
GRANT EXECUTE ON pdb_admin.insert_ligne_commande TO s1_user;

-- Sur chaque SITE : g_user n'execute QUE les procedures de routage
GRANT CREATE SESSION                     TO g_user;
GRANT EXECUTE ON pdb_admin.insertligne   TO g_user;
GRANT EXECUTE ON pdb_admin.deleteligne   TO g_user;
GRANT EXECUTE ON pdb_admin.updateligne   TO g_user;
GRANT SELECT  ON pdb_admin.LIGNECOMMANDES1 TO g_user;  -- lecture des fragments
\end{lstlisting}

Une factorisation de ces autorisations dans un rôle \texttt{site\_role} a été
préparée (et laissée en commentaire dans \texttt{04\_create\_roles.sql}) ; la
version déployée conserve des \texttt{GRANT} explicites par compte, plus
lisibles à des fins pédagogiques et auditables d'un coup d'œil.

\subsubsection*{Qui se connecte à qui}

La communication repose entièrement sur des \emph{database links}, chacun
ouvrant une session sur la base distante \textbf{au nom d'un compte de service
précis}. Le sens des flux est dissymétrique : les sites \emph{tirent} les
données de référence depuis le global, tandis que le global \emph{pousse} les
écritures vers les sites via les déclencheurs (section~\ref{sec:triggers}).

\begin{table}[h]
\centering
\begin{tabular}{@{}lllll@{}}
\toprule
\textbf{Lien} & \textbf{Créé sur} & \textbf{Cible} & \textbf{Connexion en} & \textbf{Usage} \\
\midrule
\texttt{link\_to\_global} & Site 1 & \texttt{G\_PDB}  & \texttt{s1\_user} & lecture des parents (synonymes) \\
\texttt{link\_to\_global} & Site 2 & \texttt{G\_PDB}  & \texttt{s2\_user} & lecture des parents (synonymes) \\
\texttt{link\_to\_global} & Backup & \texttt{G\_PDB}  & \texttt{b\_user}  & rechargement horaire \\
\texttt{link\_to\_site1}  & Global & \texttt{S1\_PDB} & \texttt{g\_user}  & propagation des écritures \\
\texttt{link\_to\_site2}  & Global & \texttt{S2\_PDB} & \texttt{g\_user}  & propagation des écritures \\
\bottomrule
\end{tabular}
\caption{Les cinq liens de base de données et le compte de service utilisé par
chacun.}
\label{tab:links}
\end{table}

\begin{figure}[h]
\centering
\begin{tikzpicture}[
    font=\small,
    box/.style={draw,thick,rounded corners,align=left,
                minimum width=4.0cm,minimum height=1.45cm,inner sep=6pt},
    glob/.style={box,fill=glob},
    site/.style={box,fill=site},
    bkp/.style={box,fill=bkp},
    push/.style={-{Latex[length=2.4mm]},thick},
    pull/.style={-{Latex[length=2.4mm]},thick,dashed},
]
% --- Noeuds ---
\node[glob] (g)  at (0,4.4)   {\textbf{oracle-global} (\texttt{G\_PDB})\\[2pt]
    \scriptsize\texttt{pdb\_admin} (propriétaire)\\
    \scriptsize\texttt{s1\_user}, \texttt{s2\_user}, \texttt{b\_user} (SELECT)};
\node[site] (s1) at (-5.4,1.4){\textbf{oracle-site-1} (\texttt{S1\_PDB})\\[2pt]
    \scriptsize\texttt{pdb\_admin} (propriétaire)\\
    \scriptsize\texttt{g\_user} (EXECUTE)};
\node[site] (s2) at (5.4,1.4) {\textbf{oracle-site-2} (\texttt{S2\_PDB})\\[2pt]
    \scriptsize\texttt{pdb\_admin} (propriétaire)\\
    \scriptsize\texttt{g\_user} (EXECUTE)};
\node[bkp]  (b)  at (0,-1.2)  {\textbf{oracle-backup} (\texttt{B\_PDB})\\[2pt]
    \scriptsize\texttt{pdb\_admin} (propriétaire)};

% --- Liens : deux arcs opposes par paire (pas d'etiquette sur les fleches) ---
% Trait plein  = ecritures poussees par le global (compte g_user)
% Pointilles   = lectures tirees du global (comptes s1_user / s2_user / b_user)
\draw[push] (g)  to[bend left=15] (s1);
\draw[pull] (s1) to[bend left=15] (g);
\draw[push] (g)  to[bend left=15] (s2);
\draw[pull] (s2) to[bend left=15] (g);
\draw[pull] (b) -- (g);

% --- Legende ---
\node[draw,rounded corners,fill=white,inner sep=6pt,align=left,font=\scriptsize,
      anchor=north] at (0,-2.3) {%
  \tikz{\draw[-{Latex[length=1.8mm]},thick] (0,0)--(0.7,0);}~~écritures poussées par le global (compte \texttt{g\_user})\\[3pt]
  \tikz{\draw[-{Latex[length=1.8mm]},thick,dashed] (0,0)--(0.7,0);}~~lectures tirées du global (comptes \texttt{s1\_user}, \texttt{s2\_user}, \texttt{b\_user})%
};
\end{tikzpicture}
\caption{Communication inter-nœuds. Les flèches pleines représentent les
écritures poussées par le nœud global vers les sites (via le compte
\texttt{g\_user}) ; les flèches pointillées, les lectures que les autres nœuds
tirent du global. Le détail de chaque lien et du compte qu'il emploie figure au
tableau~\ref{tab:links}.}
\label{fig:users}
\end{figure}

\subsection{Les deux scénarios de fragmentation}

Le projet implémente deux stratégies de fragmentation, déployées sur deux piles
indépendantes (répertoires \texttt{scenario-1/} et \texttt{scenario-2/}). Le
tableau~\ref{tab:scenarios} les compare ; la section~\ref{sec:fragmentation}
les détaille.

\begin{table}[h]
\centering
\begin{tabular}{@{}lll@{}}
\toprule
 & \textbf{Scénario 1 (par catégorie)} & \textbf{Scénario 2 (par volume)} \\
\midrule
Fragment R1 (Site 1) & \texttt{IDCATEG=50 $\wedge$ QTÉ$>$100} & \texttt{QTÉ $\geq$ 100} \\
Fragment R2 (Site 2) & \texttt{IDCATEG=35 $\wedge$ QTÉ$>$50}  & \texttt{QTÉ $<$ 100} \\
Disjoints            & oui & oui \\
Exhaustifs           & \textbf{non} (reste au global) & \textbf{oui} (partition complète) \\
Donnée de routage    & catégorie (jointure \texttt{PRODUITS}) & \texttt{QUANTITE} (colonne directe) \\
\bottomrule
\end{tabular}
\caption{Comparaison des deux stratégies de fragmentation horizontale.}
\label{tab:scenarios}
\end{table}

% ============================================================================
\section{Modèle de données}
% ============================================================================

Le schéma global comprend sept tables. \texttt{CLIENTS}, \texttt{FOURNISSEURS},
\texttt{EMPLOYES} et \texttt{CATEGORIES} sont des tables de référence ;
\texttt{COMMANDES}, \texttt{PRODUITS} et \texttt{LIGNECOMMANDES} portent
l'activité transactionnelle. Les relations clés sont : une \texttt{COMMANDE}
appartient à un \texttt{CLIENT} et à un \texttt{EMPLOYE} ; un \texttt{PRODUIT}
relève d'une \texttt{CATEGORIE} et d'un \texttt{FOURNISSEUR} ; une
\texttt{LIGNECOMMANDE} relie une \texttt{COMMANDE} à un \texttt{PRODUIT} avec
une quantité et une remise. Le rôle de chacune des sept tables est le suivant :

\begin{itemize}[leftmargin=1.4cm]
  \item \texttt{CLIENTS} : fiche des clients (société, contact, adresse, pays\ldots).
  \item \texttt{FOURNISSEURS} : fiche des fournisseurs approvisionnant les produits.
  \item \texttt{EMPLOYES} : personnel ayant saisi les commandes.
  \item \texttt{CATEGORIES} : nomenclature des familles de produits ; sa clé
        \texttt{IDCATEG} sert de critère de routage au scénario~1.
  \item \texttt{COMMANDES} : en-têtes de commande (client, employé, dates).
  \item \texttt{PRODUITS} : catalogue, rattaché à une catégorie et un fournisseur.
  \item \texttt{LIGNECOMMANDES} : lignes de commande (produit, quantité, remise) ;
        c'est la \textbf{seule table fragmentée} entre les sites.
\end{itemize}

Les quatre tables de référence portent les clés primaires vers lesquelles
pointent les clés étrangères des tables transactionnelles :

\begin{lstlisting}[style=sql,caption={Tables de référence du schéma global (extrait de \texttt{02\_create\_table.sql}).},label={lst:schema-ref}]
CREATE TABLE CLIENTS (
    IDCLIENT   NUMBER PRIMARY KEY,
    CODECLIENT VARCHAR2(100),
    SOCIETE    VARCHAR2(100) NOT NULL,
    CONTACT    VARCHAR2(100) NOT NULL,
    FONCTION   VARCHAR2(100) NOT NULL,
    ADRESSE    VARCHAR2(100),  VILLE  VARCHAR2(100),  REGION VARCHAR2(100),
    CP         VARCHAR2(10),   PAYS   VARCHAR2(100),
    NAISSANCE  DATE,
    TELEPHONE  VARCHAR2(100),  FAX    VARCHAR2(100)
);

CREATE TABLE FOURNISSEURS (
    IDFOUR      INTEGER PRIMARY KEY,
    SOCIETEFOUR VARCHAR2(100 CHAR),
    CONTACTFOUR VARCHAR2(100 CHAR),
    -- ... autres colonnes : fonction, adresse, ville, region, CP, telephone, fax
    PAYSFOUR    VARCHAR2(100 CHAR)
);

CREATE TABLE EMPLOYES (
    IDEMPLOYE       INTEGER PRIMARY KEY,
    NOM             VARCHAR2(100 CHAR),
    PRENOM          VARCHAR2(100 CHAR),
    FONCTIONEMPLOYE VARCHAR2(100 CHAR),
    DATEEMBAUCHE    DATE,
    -- ... autres colonnes : titre de courtoisie, naissance, adresse, ville, pays, tel
    VILLEEMPLOYE    VARCHAR2(100 CHAR)
);

CREATE TABLE CATEGORIES (
    IDCATEG        INTEGER PRIMARY KEY,
    NOMDECATEGORIE VARCHAR2(100 CHAR),
    DESCRIPTION    CLOB
);
\end{lstlisting}

Les trois tables transactionnelles, qui matérialisent ces relations par des
clés étrangères, sont créées ensuite :

\begin{lstlisting}[style=sql,caption={Tables transactionnelles du schéma global (extrait de \texttt{02\_create\_table.sql}).},label={lst:schema}]
CREATE TABLE COMMANDES (
    IDCOMMANDE    INTEGER PRIMARY KEY,
    IDEMPLOYE     INTEGER REFERENCES Employes(IDEMPLOYE),
    IDCLIENT      INTEGER REFERENCES Clients(IDCLIENT),
    DATECOMMANDE  DATE,
    DATELIVRAISON DATE, CHECK (DATELIVRAISON >= DATECOMMANDE),
    NMESSAGER     NUMBER(4,0),
    PORTNUMBER    NUMBER(4,0)
);

CREATE TABLE PRODUITS (
    IDPRODUIT    INTEGER PRIMARY KEY,
    DESIGNATION  VARCHAR2(100 CHAR),
    IDFOUR       INTEGER REFERENCES FOURNISSEURS(IDFOUR),
    IDCATEG      INTEGER REFERENCES CATEGORIES(IDCATEG),
    PRIXUNITAIRE FLOAT,
    INDISPONIBLE INTEGER, CHECK (INDISPONIBLE IN (0,1))
);

CREATE TABLE LigneCommandes (
    IDLIGNECOMMANDE INTEGER PRIMARY KEY,
    IDCOMMANDE      INTEGER REFERENCES COMMANDES(IDCOMMANDE),
    IDPRODUIT       INTEGER REFERENCES PRODUITS(IDPRODUIT),
    QUANTITE        INTEGER,
    REMISE          FLOAT
);
\end{lstlisting}

On notera que la \textbf{catégorie d'un produit n'est pas stockée} sur la ligne
de commande : elle n'est accessible que par jointure avec \texttt{PRODUITS}.
Cette particularité a une conséquence directe sur le scénario~1, où le routage
dépend de la catégorie (voir section~\ref{sec:triggers}).

% ============================================================================
\section{Fragmentation horizontale}
\label{sec:fragmentation}
% ============================================================================

\subsection{Principe}

La \emph{fragmentation horizontale} consiste à partitionner une table en
sous-ensembles de \emph{lignes} (et non de colonnes) au moyen d'un prédicat de
sélection. En algèbre relationnelle, chaque fragment $R_i$ est défini par une
sélection $\sigma_{P_i}(R)$. Deux propriétés sont recherchées :

\begin{itemize}[leftmargin=1.4cm]
  \item \textbf{Disjonction} : $\sigma_{P_i}(R) \cap \sigma_{P_j}(R) =
    \varnothing$ pour $i \neq j$, aucune ligne n'est dupliquée entre deux
    sites.
  \item \textbf{Complétude (ou exhaustivité)} : $\bigcup_i \sigma_{P_i}(R) = R$,
    aucune ligne n'est perdue.
\end{itemize}

\subsection{Scénario 1 : fragmentation par catégorie}

\begin{align*}
R_1 &= \sigma_{\,\text{IDCATEG}=50 \,\wedge\, \text{QUANTITE}>100}(\text{LigneCommandes}), \\
R_2 &= \sigma_{\,\text{IDCATEG}=35 \,\wedge\, \text{QUANTITE}>50}(\text{LigneCommandes}).
\end{align*}

Ces deux fragments sont \textbf{disjoints} (les catégories 50 et 35 sont
exclusives) mais \textbf{non exhaustifs} : toutes les lignes dont la catégorie
n'est ni 50 ni 35, ou dont la quantité est trop faible, n'appartiennent à aucun
site et ne subsistent que sur le nœud global.

\subsection{Scénario 2 : fragmentation par volume}

\begin{align*}
R_1 &= \sigma_{\,\text{QUANTITE} \geq 100}(\text{LigneCommandes}), \\
R_2 &= \sigma_{\,\text{QUANTITE} < 100}(\text{LigneCommandes}).
\end{align*}

Cette partition est à la fois disjointe \emph{et} exhaustive : toute ligne
appartient à exactement un site. Le routage ne nécessite alors aucune jointure,
puisque \texttt{QUANTITE} est une colonne directe de la table.

\subsection{Construction des fragments et intégrité référentielle}

Chaque site construit ses tables de fragment par \texttt{CREATE TABLE AS
SELECT} en filtrant les données rapatriées du global via les synonymes. Comme
cette syntaxe \textbf{n'hérite pas des contraintes} (clés primaires, clés
étrangères), celles-ci sont \textbf{reconstruites explicitement}. De plus, pour
préserver l'intégrité référentielle, le site n'importe pas seulement les lignes
de commande : il importe aussi les \texttt{PRODUITS}, \texttt{COMMANDES} et
\texttt{CLIENTS} \emph{associés} (réponse à Q3).

\begin{lstlisting}[style=sql,caption={Construction du fragment R1 sur le site~1 (extrait de \texttt{03\_create\_tables.sql}).},label={lst:fragment}]
-- Sous-ensemble des lignes de commande satisfaisant R1
CREATE TABLE LIGNECOMMANDES1 AS
    SELECT lc.*
    FROM   LIGNECOMMANDES lc
    JOIN   PRODUITS p ON lc.IDPRODUIT = p.IDPRODUIT
    WHERE  lc.QUANTITE > 100
    AND    p.IDCATEG   = 50;

-- La cle primaire n'est PAS heritee du CREATE TABLE AS SELECT
ALTER TABLE LIGNECOMMANDES1
    ADD CONSTRAINT lignecommandes1_pk PRIMARY KEY (IDLIGNECOMMANDE);

-- On ne rapatrie que les produits/commandes/clients effectivement references
CREATE TABLE PRODUITS1 AS
    SELECT DISTINCT p.*
    FROM   PRODUITS p
    JOIN   LIGNECOMMANDES1 lc ON p.IDPRODUIT = lc.IDPRODUIT;

-- ... puis on retablit les cles etrangeres entre fragments locaux
ALTER TABLE LIGNECOMMANDES1
    ADD CONSTRAINT lignecommandes1_produits1_fk
    FOREIGN KEY (IDPRODUIT) REFERENCES PRODUITS1(IDPRODUIT);
\end{lstlisting}

% ============================================================================
\section{Procédures stockées de chaque site (Q2 et Q3)}
% ============================================================================

Chaque site expose trois procédures PL/SQL (\texttt{insertligne},
\texttt{deleteligne} et \texttt{updateligne}) qui encapsulent toute la
logique de manipulation de son fragment. Elles sont appelées à distance par les
déclencheurs du nœud global (section~\ref{sec:triggers}).

\subsection{Insertion avec rapatriement des parents}

La procédure \texttt{insertligne} ne se contente pas d'insérer la ligne : elle
\textbf{garantit au préalable la présence des lignes parentes}. Si le produit,
la commande ou le client correspondant est absent du fragment, il est importé à
la volée depuis le global. C'est cette discipline qui maintient l'intégrité
référentielle malgré la fragmentation.

\begin{lstlisting}[style=sql,caption={Insertion avec rapatriement automatique des parents (extrait, site~1).},label={lst:insert}]
CREATE OR REPLACE PROCEDURE insertligne (
    p_idlignecommande IN LIGNECOMMANDES1.IDLIGNECOMMANDE%TYPE,
    p_idcommande      IN LIGNECOMMANDES1.IDCOMMANDE%TYPE,
    p_idproduit       IN LIGNECOMMANDES1.IDPRODUIT%TYPE,
    p_quantite        IN LIGNECOMMANDES1.QUANTITE%TYPE,
    p_remise          IN LIGNECOMMANDES1.REMISE%TYPE
) IS
    v_count    INTEGER;
    v_idclient COMMANDES1.IDCLIENT%TYPE;
BEGIN
    -- Le produit existe-t-il deja localement ? Sinon, on l'importe du global.
    SELECT COUNT(*) INTO v_count FROM PRODUITS1 WHERE IDPRODUIT = p_idproduit;
    IF v_count = 0 THEN
        INSERT INTO PRODUITS1 SELECT * FROM PRODUITS WHERE IDPRODUIT = p_idproduit;
    END IF;

    -- Idem pour la commande (et le client qu'elle reference)
    SELECT COUNT(*) INTO v_count FROM COMMANDES1 WHERE IDCOMMANDE = p_idcommande;
    IF v_count = 0 THEN
        SELECT IDCLIENT INTO v_idclient FROM COMMANDES WHERE IDCOMMANDE = p_idcommande;
        SELECT COUNT(*) INTO v_count FROM CLIENTS1 WHERE IDCLIENT = v_idclient;
        IF v_count = 0 THEN
            INSERT INTO CLIENTS1 SELECT * FROM CLIENTS WHERE IDCLIENT = v_idclient;
        END IF;
        INSERT INTO COMMANDES1 SELECT * FROM COMMANDES WHERE IDCOMMANDE = p_idcommande;
    END IF;

    INSERT INTO LIGNECOMMANDES1 (IDLIGNECOMMANDE, IDCOMMANDE, IDPRODUIT, QUANTITE, REMISE)
    VALUES (p_idlignecommande, p_idcommande, p_idproduit, p_quantite, p_remise);
    COMMIT;
EXCEPTION
    WHEN OTHERS THEN ROLLBACK; RAISE;
END insertligne;
/
\end{lstlisting}

\subsection{Suppression en cascade « ascendante »}

Symétriquement, \texttt{deleteligne} supprime la ligne puis \textbf{nettoie les
parents devenus orphelins} : si la commande n'a plus aucune ligne, elle est
supprimée ; si le client n'a plus aucune commande, il l'est à son tour. Les
produits, en revanche, sont des données de référence partagées et ne sont
jamais supprimés.

\begin{lstlisting}[style=sql,caption={Suppression et nettoyage des parents orphelins (extrait).},label={lst:delete}]
DELETE FROM LIGNECOMMANDES1 WHERE IDLIGNECOMMANDE = p_idlignecommande;

-- La commande a-t-elle encore des lignes ?
SELECT COUNT(*) INTO v_count FROM LIGNECOMMANDES1 WHERE IDCOMMANDE = v_idcommande;
IF v_count = 0 THEN
    DELETE FROM COMMANDES1 WHERE IDCOMMANDE = v_idcommande;
    -- Le client a-t-il encore des commandes ?
    SELECT COUNT(*) INTO v_count FROM COMMANDES1 WHERE IDCLIENT = v_idclient;
    IF v_count = 0 THEN
        DELETE FROM CLIENTS1 WHERE IDCLIENT = v_idclient;
    END IF;
END IF;
\end{lstlisting}

\subsection{Habilitations}

Les procédures sont exécutées au travers du lien de base de données par un
compte de service dédié, \texttt{g\_user}. Celui-ci reçoit donc les droits
\texttt{EXECUTE} sur les trois procédures, \texttt{SELECT} sur les tables de
fragment, ainsi que \texttt{CREATE SESSION}, et rien d'autre. Ce compte est une
illustration directe du principe du moindre privilège détaillé en
section~\ref{sec:users}.

% ============================================================================
\section{Synchronisation par déclencheurs (Q4)}
\label{sec:triggers}
% ============================================================================

C'est le cœur du système. Trois déclencheurs (\emph{triggers}) posés sur la
table \texttt{LIGNECOMMANDES} du nœud global interceptent chaque \texttt{INSERT},
\texttt{DELETE} et \texttt{UPDATE} et propagent la modification vers le site
propriétaire, en appelant à distance les procédures de la
section précédente. L'application n'écrit donc \textbf{que sur le global} : le
routage est entièrement automatique et transparent.

\subsection{Routage à l'insertion}

Le déclencheur d'insertion résout d'abord la catégorie du produit (absente de
la ligne, comme noté plus haut), puis dirige la ligne vers Site~1 et/ou Site~2
selon les prédicats $R_1$ et $R_2$.

\begin{lstlisting}[style=sql,caption={Déclencheur de routage à l'insertion, scénario~1 (extrait de \texttt{06\_create\_triggers.sql}).},label={lst:trig-insert}]
CREATE OR REPLACE TRIGGER SYC_INSERT_LIGNE
AFTER INSERT ON pdb_admin.LIGNECOMMANDES
FOR EACH ROW
DECLARE
    PRAGMA AUTONOMOUS_TRANSACTION;          -- (voir explications ci-dessous)
    v_idcateg PRODUITS.IDCATEG%TYPE;
BEGIN
    -- La categorie n'est pas sur la ligne : on la lit dans PRODUITS
    SELECT IDCATEG INTO v_idcateg FROM PRODUITS WHERE IDPRODUIT = :NEW.IDPRODUIT;

    -- Route vers Site 1 si la ligne satisfait R1
    IF v_idcateg = 50 AND :NEW.QUANTITE > 100 THEN
        EXECUTE IMMEDIATE 'BEGIN pdb_admin.insertligne@link_to_site1(:1,:2,:3,:4,:5); END;'
            USING :NEW.IDLIGNECOMMANDE, :NEW.IDCOMMANDE, :NEW.IDPRODUIT,
                  :NEW.QUANTITE, :NEW.REMISE;
    END IF;

    -- Route vers Site 2 si la ligne satisfait R2
    IF v_idcateg = 35 AND :NEW.QUANTITE > 50 THEN
        EXECUTE IMMEDIATE 'BEGIN pdb_admin.insertligne@link_to_site2(:1,:2,:3,:4,:5); END;'
            USING :NEW.IDLIGNECOMMANDE, :NEW.IDCOMMANDE, :NEW.IDPRODUIT,
                  :NEW.QUANTITE, :NEW.REMISE;
    END IF;
EXCEPTION
    WHEN NO_DATA_FOUND THEN NULL;           -- produit absent => aucun fragment
    WHEN OTHERS THEN
        RAISE_APPLICATION_ERROR(-20010, 'SYC_INSERT_LIGNE: ' || SQLERRM);
END SYC_INSERT_LIGNE;
/
\end{lstlisting}

\subsection{Deux décisions de conception déterminantes}

Le code ci-dessus comporte deux choix techniques qui méritent d'être justifiés,
car ils contournent deux erreurs Oracle bien précises.

\paragraph{1. \texttt{EXECUTE IMMEDIATE} plutôt qu'un appel direct.}
Oracle~23c valide les références aux procédures distantes \emph{à la compilation}
du déclencheur. Or les conteneurs des sites démarrent \emph{après} le global :
au moment où le déclencheur est compilé, les procédures distantes n'existent pas
encore, et un appel direct échouerait avec l'erreur \texttt{PLS-00352}. En
passant par du SQL dynamique (\texttt{EXECUTE IMMEDIATE}), la résolution est
\textbf{différée à l'exécution}, ce qui rompt cette dépendance de compilation.

\paragraph{2. \texttt{PRAGMA AUTONOMOUS\_TRANSACTION}.}
Les procédures distantes effectuent un \texttt{COMMIT}. Sans cette directive, ce
\texttt{COMMIT} se produirait au sein de la transaction distribuée pilotée par
le déclencheur, ce qui déclencherait l'erreur \texttt{ORA-02089} (« \texttt{COMMIT}
interdit dans une session subordonnée »). La directive isole chaque appel de
site dans sa \textbf{propre transaction autonome}.

\medskip
\noindent\textbf{Contrepartie assumée.} Cette autonomie transactionnelle a un
prix : si la transaction globale est ensuite annulée (\texttt{ROLLBACK}), le
site a déjà validé sa copie de la ligne. Il n'y a donc \textbf{pas d'atomicité
distribuée}. Ce compromis est acceptable dans un cadre pédagogique ; un système
de production exigerait une validation en deux phases (2PC/XA) ou une file de
messages fiable (Oracle Advanced Queuing). Ce point est repris en
section~\ref{sec:limites}.

\subsection{Le cas subtil de la mise à jour}

Une mise à jour peut faire \emph{entrer} une ligne dans un fragment, l'en faire
\emph{sortir}, ou l'y \emph{maintenir}. Le déclencheur de mise à jour compare
donc l'appartenance au fragment \emph{avant} et \emph{après} la modification, et
choisit l'action adéquate parmi quatre cas, pour chaque site
(tableau~\ref{tab:update}).

\begin{table}[h]
\centering
\begin{tabular}{@{}lll@{}}
\toprule
\textbf{Avant} & \textbf{Après} & \textbf{Action sur le site} \\
\midrule
dans le fragment & dans le fragment & \texttt{updateligne} (mise à jour sur place) \\
dans le fragment & hors fragment    & \texttt{deleteligne} (la ligne quitte le site) \\
hors fragment    & dans le fragment & \texttt{insertligne} (la ligne entre sur le site) \\
hors fragment    & hors fragment    & aucune action \\
\bottomrule
\end{tabular}
\caption{Les quatre cas gérés par \texttt{SYC\_UPDATE\_LIGNE}, par site.}
\label{tab:update}
\end{table}

\begin{lstlisting}[style=sql,caption={Logique de transition de fragment à la mise à jour (extrait, partie Site~1).},label={lst:trig-update}]
v_in_s1_old := (v_old_categ = 50 AND :OLD.QUANTITE > 100);
v_in_s1_new := (v_new_categ = 50 AND :NEW.QUANTITE > 100);

IF v_in_s1_old AND v_in_s1_new THEN
    EXECUTE IMMEDIATE 'BEGIN pdb_admin.updateligne@link_to_site1(:1,:2,:3,:4); END;'
        USING :NEW.IDLIGNECOMMANDE, :NEW.IDPRODUIT, :NEW.QUANTITE, :NEW.REMISE;
ELSIF v_in_s1_old AND NOT v_in_s1_new THEN
    EXECUTE IMMEDIATE 'BEGIN pdb_admin.deleteligne@link_to_site1(:1); END;'
        USING :OLD.IDLIGNECOMMANDE;
ELSIF NOT v_in_s1_old AND v_in_s1_new THEN
    EXECUTE IMMEDIATE 'BEGIN pdb_admin.insertligne@link_to_site1(:1,:2,:3,:4,:5); END;'
        USING :NEW.IDLIGNECOMMANDE, :NEW.IDCOMMANDE, :NEW.IDPRODUIT,
              :NEW.QUANTITE, :NEW.REMISE;
END IF;
\end{lstlisting}

Dans le scénario~2, cette logique se simplifie en une comparaison du seul
franchissement du seuil \texttt{QUANTITE = 100}, puisque le routage ne dépend
plus de la catégorie.

% ============================================================================
\section{Requêtes distribuées et optimisation (Q5 et Q6)}
% ============================================================================

\subsection{Analyse et optimisation d'une requête locale (Q5)}

La requête analytique étudiée compte le nombre de commandes par client pour
l'année 2026. Sur un schéma « froid » (sans index), le plan d'exécution révèle
deux \texttt{TABLE ACCESS FULL} (balayages complets) et une jointure par hachage
coûteuse. Deux leviers d'optimisation sont appliqués :

\begin{enumerate}[leftmargin=1.4cm]
  \item \textbf{Réécriture du prédicat de date.} La forme
    \texttt{EXTRACT(YEAR FROM DATECOMMANDE)=2026} \emph{empêche} l'usage d'un
    index B-tree (la valeur indexée n'est pas la valeur calculée). On la
    remplace par un prédicat d'intervalle
    \texttt{DATECOMMANDE >= DATE '2026-01-01' AND < DATE '2027-01-01'}, qui
    autorise un \texttt{INDEX RANGE SCAN}.
  \item \textbf{Création d'index} sur \texttt{COMMANDES(DATECOMMANDE)} et
    \texttt{COMMANDES(IDCLIENT)}, ce qui permet à l'optimiseur de remplacer le
    balayage complet par un parcours d'index et la jointure par hachage par une
    jointure par boucles imbriquées (\texttt{NESTED LOOPS}).
\end{enumerate}

\begin{lstlisting}[style=sql,caption={Capture d'un plan d'exécution et création d'index (extrait de \texttt{07\_query\_analysis.sql}).},label={lst:explain}]
-- Plan AVANT optimisation
EXPLAIN PLAN SET STATEMENT_ID = 'before_index' FOR
SELECT c.IDCLIENT, c.SOCIETE, COUNT(cmd.IDCOMMANDE) AS NB_COMMANDES
FROM   CLIENTS c JOIN COMMANDES cmd ON c.IDCLIENT = cmd.IDCLIENT
WHERE  EXTRACT(YEAR FROM cmd.DATECOMMANDE) = 2026     -- bloque l'index
GROUP  BY c.IDCLIENT, c.SOCIETE
ORDER  BY NB_COMMANDES DESC;

SELECT * FROM TABLE(DBMS_XPLAN.DISPLAY(
    filter_preds => 'STATEMENT_ID = ''before_index'''));

-- Index correctifs
CREATE INDEX idx_commandes_date     ON COMMANDES(DATECOMMANDE);
CREATE INDEX idx_commandes_idclient ON COMMANDES(IDCLIENT);

-- Requete optimisee : predicat d'intervalle => INDEX RANGE SCAN
EXPLAIN PLAN SET STATEMENT_ID = 'after_index' FOR
SELECT c.IDCLIENT, c.SOCIETE, COUNT(cmd.IDCOMMANDE) AS NB_COMMANDES
FROM   CLIENTS c JOIN COMMANDES cmd ON c.IDCLIENT = cmd.IDCLIENT
WHERE  cmd.DATECOMMANDE >= DATE '2026-01-01'
AND    cmd.DATECOMMANDE <  DATE '2027-01-01'
GROUP  BY c.IDCLIENT, c.SOCIETE
ORDER  BY NB_COMMANDES DESC;
\end{lstlisting}

\subsection{Reconstruction d'une vue globale (Q6)}

Pour calculer le chiffre d'affaires 2026 par catégorie, on recombine les deux
fragments distants au travers des liens de base de données. Le revenu d'une
ligne vaut $\text{QUANTITE} \times \text{PRIXUNITAIRE} \times (1 -
\text{REMISE})$. On emploie \texttt{UNION ALL} (et non \texttt{UNION}) : les
fragments étant disjoints par construction, aucune ligne n'apparaît deux fois,
et \texttt{UNION ALL} évite la passe de déduplication superflue.

\begin{lstlisting}[style=sql,caption={Chiffre d'affaires distribué reconstruit via les liens de base de données (extrait de \texttt{08\_distributed\_query.sql}).},label={lst:distrib}]
SELECT IDCATEG, ROUND(SUM(CA_LIGNE), 2) AS CHIFFRE_AFFAIRES_2026
FROM (
    -- Contribution du Site 1
    SELECT p.IDCATEG,
           lc.QUANTITE * p.PRIXUNITAIRE * (1 - NVL(lc.REMISE, 0)) AS CA_LIGNE
    FROM   pdb_admin.LIGNECOMMANDES1@link_to_site1 lc
    JOIN   pdb_admin.PRODUITS1@link_to_site1  p ON lc.IDPRODUIT  = p.IDPRODUIT
    JOIN   pdb_admin.COMMANDES1@link_to_site1 c ON lc.IDCOMMANDE = c.IDCOMMANDE
    WHERE  c.DATECOMMANDE >= DATE '2026-01-01'
    AND    c.DATECOMMANDE <  DATE '2027-01-01'

    UNION ALL          -- fragments disjoints : pas de doublon possible

    -- Contribution du Site 2
    SELECT p.IDCATEG,
           lc.QUANTITE * p.PRIXUNITAIRE * (1 - NVL(lc.REMISE, 0)) AS CA_LIGNE
    FROM   pdb_admin.LIGNECOMMANDES2@link_to_site2 lc
    JOIN   pdb_admin.PRODUITS2@link_to_site2  p ON lc.IDPRODUIT  = p.IDPRODUIT
    JOIN   pdb_admin.COMMANDES2@link_to_site2 c ON lc.IDCOMMANDE = c.IDCOMMANDE
    WHERE  c.DATECOMMANDE >= DATE '2026-01-01'
    AND    c.DATECOMMANDE <  DATE '2027-01-01'
)
GROUP BY IDCATEG
ORDER BY CHIFFRE_AFFAIRES_2026 DESC;
\end{lstlisting}

% ============================================================================
\section{Nœud de sauvegarde}
% ============================================================================

Un quatrième nœud assure une \textbf{copie de secours} du jeu de données global.
Il héberge un schéma miroir (sans clés étrangères, volontairement) et une
procédure \texttt{sync\_from\_global} qui effectue un \textbf{rechargement
complet} : vidage de toutes les tables (\texttt{TRUNCATE}, des feuilles vers la
racine) puis réinsertion intégrale depuis le global (de la racine vers les
feuilles, pour respecter l'ordre des dépendances). Cette procédure est
planifiée toutes les heures par \texttt{DBMS\_SCHEDULER}.

\begin{lstlisting}[style=sql,caption={Procédure de resynchronisation et job horaire (extrait de \texttt{04\_create\_sync\_job.sql}).},label={lst:backup}]
CREATE OR REPLACE PROCEDURE pdb_admin.sync_from_global AS
BEGIN
    -- Vidage des enfants vers les parents
    EXECUTE IMMEDIATE 'TRUNCATE TABLE pdb_admin.LIGNECOMMANDES';
    EXECUTE IMMEDIATE 'TRUNCATE TABLE pdb_admin.COMMANDES';
    -- ... (autres tables)

    -- Rechargement des parents vers les enfants
    INSERT INTO pdb_admin.CATEGORIES SELECT * FROM pdb_admin.CATEGORIES@link_to_global;
    -- ...
    INSERT INTO pdb_admin.LIGNECOMMANDES SELECT * FROM pdb_admin.LIGNECOMMANDES@link_to_global;
    COMMIT;
END sync_from_global;
/

BEGIN
  DBMS_SCHEDULER.CREATE_JOB(
    job_name        => 'BACKUP_SYNC_JOB',
    job_type        => 'STORED_PROCEDURE',
    job_action      => 'pdb_admin.sync_from_global',
    start_date      => SYSTIMESTAMP,
    repeat_interval => 'FREQ=HOURLY;INTERVAL=1',
    enabled         => TRUE);
END;
/
\end{lstlisting}

% ============================================================================
\section{Tableau de bord web}
% ============================================================================

Une application \textbf{Next.js} (App Router, React~19) offre une interface
d'observation et de test. Côté serveur, des \emph{routes API} se connectent à
Oracle via le pilote \texttt{node-oracledb}, en gérant un \textbf{pool de
connexions} par nœud. Le pool est mémorisé dans \texttt{globalThis} pour
survivre aux rechargements à chaud du serveur de développement.

\begin{lstlisting}[style=ts,caption={Gestion d'un pool de connexions par nœud (extrait de \texttt{lib/oracle.ts}).},label={lst:pool}]
async function ensurePool(alias: NodeAlias) {
  if (g._oraPools!.has(alias)) return;
  try {
    await oracledb.createPool({
      poolAlias: alias,
      ...configs[alias](),
      poolMin: 0, poolMax: 4, poolIncrement: 1,
      poolTimeout: 60,
      connectTimeout: 5, // secondes ; evite les fuites de connexion
    });
    g._oraPools!.add(alias);
  } catch (err: unknown) {
    const msg = err instanceof Error ? err.message : String(err);
    if (msg.includes('NJS-046') || msg.includes('already exists')) {
      g._oraPools!.add(alias);     // pool deja cree par un hot-reload
    } else { throw err; }
  }
}
\end{lstlisting}

La route d'insertion illustre bien la \emph{transparence} du dispositif :
l'application \textbf{insère uniquement sur le global}, attend brièvement la
propagation par les déclencheurs, puis interroge les deux sites pour
\emph{constater} où la ligne a été routée et le restituer à l'interface.

\begin{lstlisting}[style=ts,caption={Insertion sur le global puis observation du routage (extrait de \texttt{app/api/insert/route.ts}).},label={lst:insert-api}]
// 1) Insertion sur le noeud global : les triggers se declenchent seuls
await run('global',
  `INSERT INTO LIGNECOMMANDES (IDLIGNECOMMANDE, IDCOMMANDE, IDPRODUIT, QUANTITE, REMISE)
   VALUES (:1, :2, :3, :4, :5)`,
  [idlignecommande, idcommande, idproduit, quantite, remise]);

// 2) On laisse aux triggers le temps de propager
await sleep(700);

// 3) On constate sur quel(s) site(s) la ligne a atterri
const [s1Rows, s2Rows] = await Promise.allSettled([
  query('site1', findSql('LIGNECOMMANDES1'), [idlignecommande]),
  query('site2', findSql('LIGNECOMMANDES2'), [idlignecommande]),
]);
\end{lstlisting}

Les autres routes exposent l'état de santé des nœuds (\texttt{/api/health}, qui
vérifie aussi l'état des déclencheurs et la dernière exécution du job de
sauvegarde) et le contenu des fragments (\texttt{/api/fragments}). La page
React rafraîchit l'état de santé toutes les dix secondes.

% ============================================================================
\section{Démonstration et captures d'écran}
\label{sec:captures}
% ============================================================================

Cette section regroupe les captures d'écran illustrant le système en
fonctionnement : l'orchestration des conteneurs, l'interface du tableau de bord,
puis le déroulement des tests de manipulation (insertion, suppression, mise à
jour) qui valident la synchronisation par déclencheurs
(section~\ref{sec:triggers}).

\subsection{Conteneurs Docker en cours d'exécution}

\begin{figure}[h]
\centering
\screenshot{captures/docker-ps.png}{Sortie de \texttt{docker compose ps} montrant les quatre nœuds (global, site~1, site~2, sauvegarde) à l'état \texttt{healthy}.}
\caption{Les conteneurs Oracle orchestrés par Docker Compose, tous démarrés.}
\label{fig:cap-docker}
\end{figure}

\subsection{Tableau de bord web}

\begin{figure}[h]
\centering
\screenshot{captures/dashboard.png}{Interface Next.js : formulaire d'insertion, badges de nœud et panneau des fragments.}
\caption{Le tableau de bord affichant l'état des fragments répartis.}
\label{fig:cap-dashboard}
\end{figure}

\subsection{Tests de manipulation des données}

Les trois captures suivantes documentent un aller-retour complet sur la table
\texttt{LIGNECOMMANDES} du nœud global et la propagation attendue vers le site
propriétaire.

\begin{figure}[h]
\centering
\screenshot{captures/test-insert.png}{Insertion d'une ligne sur le global, puis vérification de sa présence dans le fragment du site routé.}
\caption{Test d'\textbf{insertion} et propagation vers le site.}
\label{fig:cap-insert}
\end{figure}

\begin{figure}[h]
\centering
\screenshot{captures/test-delete.png}{Suppression d'une ligne sur le global et disparition de sa copie sur le site (avec nettoyage des parents orphelins).}
\caption{Test de \textbf{suppression} et nettoyage en cascade.}
\label{fig:cap-delete}
\end{figure}

\begin{figure}[h]
\centering
\screenshot{captures/test-update.png}{Mise à jour modifiant l'appartenance d'une ligne et provoquant sa migration d'un fragment à l'autre.}
\caption{Test de \textbf{mise à jour} (changement de fragment).}
\label{fig:cap-update}
\end{figure}

% ============================================================================
\section{Avantages de l'architecture}
% ============================================================================

\begin{description}[leftmargin=0.6cm,style=nextline]
  \item[\textbf{Transparence de localisation}] L'application n'a qu'un seul
    point d'écriture (le global) et ignore où résident physiquement les
    données. Synonymes et déclencheurs masquent entièrement la répartition.
  \item[\textbf{Reproductibilité totale (infrastructure as code)}] L'intégralité
    du système (topologie, schémas, données, logique) est décrite par des
    fichiers versionnés. Un \texttt{docker compose up} reconstruit
    l'environnement à l'identique ; les scripts \texttt{purge.sh} et
    \texttt{rsite-1.sh} facilitent la remise à zéro.
  \item[\textbf{Scripts idempotents}] L'emploi systématique de
    \texttt{CREATE OR REPLACE} et de blocs \texttt{DROP \ldots\ EXCEPTION WHEN
    OTHERS} rend les scripts ré-exécutables sans erreur.
  \item[\textbf{Intégrité référentielle préservée par fragment}] Le
    rapatriement automatique des lignes parentes et la reconstruction des clés
    garantissent que chaque site reste un sous-schéma cohérent et interrogeable
    de façon autonome.
  \item[\textbf{Découplage à la compilation}] Le recours à
    \texttt{EXECUTE IMMEDIATE} affranchit le système de l'ordre de démarrage des
    conteneurs, point notoirement délicat des déploiements distribués.
  \item[\textbf{Optimisation mesurable}] La démarche
    \texttt{EXPLAIN PLAN} avant/après objective le gain apporté par les index et
    la réécriture des prédicats.
  \item[\textbf{Reconstruction efficace}] La disjonction garantie des fragments
    autorise un \texttt{UNION ALL} sans déduplication pour reconstituer la vue
    globale.
  \item[\textbf{Observabilité}] Le tableau de bord et le job de sauvegarde
    horaire offrent un suivi opérationnel immédiat de l'état du système.
\end{description}

% ============================================================================
\section{Inconvénients et limites}
\label{sec:limites}
% ============================================================================

\begin{description}[leftmargin=0.6cm,style=nextline]
  \item[\textbf{Absence d'atomicité distribuée}] C'est la limite majeure. La
    directive \texttt{PRAGMA AUTONOMOUS\_TRANSACTION} fait que les écritures de
    site sont validées indépendamment de la transaction globale : un
    \texttt{ROLLBACK} sur le global \emph{ne défait pas} une ligne déjà inscrite
    sur un site. En cas d'échec partiel, les fragments peuvent diverger du
    global. Une vraie cohérence exigerait une validation en deux phases (2PC/XA)
    ou une file de messages transactionnelle (Oracle AQ).
  \item[\textbf{Point de défaillance unique (SPOF)}] Tout le routage repose sur
    le nœud global et ses déclencheurs. Si le global est indisponible, plus
    aucune écriture ne peut être propagée ; les sites ne sont pas autonomes en
    écriture.
  \item[\textbf{Sauvegarde grossière (RPO élevé)}] Le nœud de secours procède
    par rechargement complet (\texttt{TRUNCATE} + réinsertion) toutes les
    heures : la perte de données potentielle peut atteindre une heure, le coût
    est proportionnel au volume total (et non aux changements), et le
    \texttt{TRUNCATE} (DDL auto-validé) laisse transitoirement les tables vides,
    sans cohérence de lecture pendant la resynchronisation.
  \item[\textbf{Sécurité perfectible}] Les mots de passe figurent en clair dans
    les scripts et le fichier Compose. Les habilitations sont accordées par
    \texttt{GRANT} directs plutôt que par rôles (le script de rôles est
    d'ailleurs laissé en commentaire), et un compte de service unique
    (\texttt{g\_user}) est partagé : le principe de moindre privilège n'est pas
    respecté.
  \item[\textbf{Génération de clés primaires fragile}] Le tableau de bord
    calcule la prochaine clé par \texttt{MAX(IDLIGNECOMMANDE)+1}, ce qui est
    sujet aux conditions de course en cas d'insertions concurrentes ; une
    séquence Oracle serait préférable.
  \item[\textbf{Propagation observée par temporisation}] La route d'insertion
    attend 700~ms « en aveugle » avant de constater le routage : ce délai
    arbitraire peut s'avérer trop court (faux négatif) ou inutilement long.
  \item[\textbf{Coût des jointures distribuées}] Les requêtes Q6 rapatrient les
    lignes des sites au travers des liens de base de données ; sur de gros
    volumes, ces aller-retours réseau deviennent un goulot d'étranglement.
  \item[\textbf{Asymétrie du scénario~1}] Sa fragmentation n'étant pas
    exhaustive, une part des lignes n'existe que sur le global, ce qui complique
    le raisonnement (une insertion peut n'aboutir sur aucun site).
\end{description}

% ============================================================================
\section{Pistes d'amélioration}
% ============================================================================

\begin{enumerate}[leftmargin=1.4cm]
  \item \textbf{Cohérence forte} : remplacer les transactions autonomes par une
    validation en deux phases, ou adopter une réplication par
    \emph{vues matérialisées} avec rafraîchissement à la validation
    (\texttt{REFRESH FAST ON COMMIT}) ou par Oracle GoldenGate.
  \item \textbf{Séquences Oracle} pour la génération des identifiants, en lieu et
    place de \texttt{MAX+1}.
  \item \textbf{Sécurité} : rôles applicatifs à privilèges minimaux, secrets
    gérés hors du dépôt (Docker secrets, coffre-fort), comptes de service
    distincts par usage.
  \item \textbf{Sauvegarde incrémentale} fondée sur la capture de changements
    (\emph{materialized view logs}, CDC) ou sur RMAN, afin de réduire le RPO et
    le coût.
  \item \textbf{Confirmation de propagation} fondée sur un événement ou un
    sondage borné plutôt que sur un délai fixe dans le tableau de bord.
  \item \textbf{Autonomie des sites} : autoriser des écritures locales avec
    réconciliation, pour atténuer la dépendance au nœud global.
\end{enumerate}

% ============================================================================
\section{Conclusion}
% ============================================================================

Ce projet met en œuvre, de bout en bout, une chaîne complète de gestion de
données réparties : fragmentation horizontale paramétrable (par catégorie ou par
volume), procédures stockées garantissant l'intégrité de chaque fragment,
synchronisation automatique par déclencheurs avec routage transparent,
optimisation de requêtes objectivée par les plans d'exécution, reconstruction
distribuée du chiffre d'affaires, sauvegarde planifiée et tableau de bord
d'observation. L'ensemble est entièrement reproductible grâce à Docker Compose
et à des scripts idempotents.

Les choix techniques les plus instructifs (le SQL dynamique pour briser la
dépendance de compilation, les transactions autonomes pour permettre la
validation distante) illustrent bien les arbitrages propres aux systèmes
distribués : ici, la \emph{disponibilité} et la \emph{simplicité} de mise en
œuvre ont été privilégiées au détriment de la \emph{cohérence forte}. Cette
tension, caractéristique des bases de données réparties, constitue précisément
le principal enseignement du projet, et trace la feuille de route d'une
éventuelle version de production : validation en deux phases, sécurité durcie et
sauvegarde incrémentale.

\end{document}
